\section{Robustness to Frequently-Executed Tasks Sample Restriction}
\label{app:frequency_robustness}

The task sequences underlying our empirical analysis combine, for each occupation, the full set of O*NET tasks regardless of how often a task is actually performed on the job.
This raises the concern that rarely-executed tasks, which a worker might perform only a few times a year, lengthen the workflow and insert ``filler'' steps between the tasks that structure day-to-day production.
If they do, the average length of the AI chains we document in Prediction~\#1, the dispersion of AI-able steps we measure in Prediction~\#2, and the adjacency between AI-executed steps we exploit in Prediction~\#3 may partly reflect infrequent tasks rather than the operative structure of work.
This appendix re-examines all three predictions on samples restricted to the tasks that workers report performing frequently.
All three results continue to hold.

Rather than re-querying the language model to re-sequence each occupation's surviving tasks, we prune the workflows already generated for the main analysis.
Within each occupation we drop the tasks that fail the frequency filter (discussed below), keep the surviving tasks in their original order, and recompute the average AI chain length, the empirical fragmentation index, and the neighbor-execution flags on the pruned sequence before re-estimating the corresponding statistics. 
Pruning the existing sequences rather than re-prompting for a fresh ordering isolates the question of interest, namely whether infrequent tasks contaminate our measures, while holding the LLM's sequencing logic fixed.

To keep the sample definition consistent across the three predictions, we restrict each cut to the occupations whose frequency-pruned workflow retains at least five tasks, and we evaluate all three predictions on this common set of occupations.
The five-task floor is applied to the pruned workflow, that is, to the steps that survive the frequency filter.
For Predictions~\#1 and~\#2 the chain length and fragmentation index are computed over each occupation's full pruned workflow; for Prediction~\#3 the neighbor regression uses the eligible tasks of the same occupations, namely those with two neighbors on either side, which a five-task workflow always admits.\footnote{For comparability across cuts, the fragmentation regressions z-score all three variables (the AI-execution share, the empirical fragmentation index, and AI exposure) within each cut, so the reported coefficients are standardized betas: a one-standard-deviation change in fragmentation maps to a $\beta$-standard-deviation change in the AI-execution share. For the neighbor logits we report average marginal effects with analytic robust standard errors clustered at the DWA level, rather than the bootstrap of Table~\ref{tab:DWA_regression_aiExecution_mainSample}. Point estimates for the ``all tasks'' specification are identical to those reported in the main text.}

\subsection{Constructing the Frequent-Task Samples}
\label{app:frequency_robustness_construction}

We measure how often a task is performed using the O*NET \emph{Frequency of Task} (FT) module.
O*NET asks the incumbent workers it surveys in each occupation how often they perform each of the occupation's tasks, offering seven answer categories ordered from least to most frequent: \textit{``yearly or less,'' ``more than yearly,'' ``more than monthly,'' ``more than weekly,'' ``daily,'' ``several times daily,''} and \textit{``hourly or more,''} and it reports the percentage of respondents choosing each category.
A task thus comes not with a single frequency but with a distribution of workers' responses over the seven categories, so restricting the sample to frequently-performed tasks requires two choices: (1) which categories count as ``frequent,'' and (2) how large the share of workers reporting the task in those categories must be.
We keep a task when the combined share of respondents in the designated frequent categories meets a threshold.

We consider three increasingly demanding notions of ``frequent.''
\emph{Daily$+$} counts tasks performed daily, several times daily, or hourly or more; \emph{SeveralDaily$+$} counts those performed several times daily or hourly or more; and \emph{Hourly$+$} counts only those performed hourly or more.
Each notion is evaluated at four thresholds: 20\%, 35\%, 50\%, and 65\% of incumbent worker responses.
A task is retained under the ``Hourly$+$ $\geq 20\%$'' cut, for example, if at least one in five workers reports performing it hourly or more.
To see how the cuts relate to one another, consider a task that 60\% of surveyed workers report performing several times daily, with the remaining responses spread over less frequent categories.
The task survives the SeveralDaily$+$ cuts at the 20\%, 35\%, and 50\% thresholds, and because responses in a stricter category also count toward every looser logic, it survives the corresponding Daily$+$ cuts as well.
Nothing stricter retains it, in either direction: no worker reports the task hourly or more, so it fails every Hourly$+$ cut, and its frequent share falls short of 65\%, so it fails the strictest threshold under both remaining logics.
Each logic-threshold pair therefore defines its own subsample of tasks and occupations.
We estimate every prediction separately on each subsample and present the results for all twelve cuts together in a grid, three logics by four thresholds, for ease of comparison.
We take the three predictions in turn.

\subsection{Prediction \#1: The Average AI Chain Length}
\label{app:frequency_robustness_chains}

Prediction~\#1 concerns the formation of AI chains.
We measure it, as in Section~\ref{sec:chainLength_prediction}, by the average length of a maximal run of contiguous AI-executed steps, where a step is AI-executed if it is labeled augmented or automated, pooled across all chains in the sample.
A longer-than-random average chain could, however, be delivered mechanically by the composition of each pruned sample.
We therefore benchmark every cut against a within-occupation task position-reshuffle placebo: we redraw $1{,}000$ reshuffles of task positions within occupations, re-apply the frequency filter, and recompute the statistic, producing a null distribution that holds the composition of the sample fixed while destroying the workflow ordering.
Figure~\ref{fig:chainlength_frequency} reports the observed average against this null as a forest plot: each row is a frequency cut, with the ``all tasks'' baseline analysis in the main manuscript on top; the grey bar spans the 10th to 90th percentiles of the reshuffle null, the vertical tick marks the null mean, and the dot is the observed estimate, red when it falls outside the null band and blue when inside, with the observed percentile printed in the right margin.
The top all-tasks row reproduces the full-sample value of $1.45$.
Chains shorten modestly as the sample is restricted to frequent tasks, to between $1.24$ and $1.31$ across most cuts, ticking back up only in the two sparsest Hourly$+$ corners ($1.40$ and $1.50$, with $75$ and $20$ occupations).
In the full sample the observed value sits at the 100th percentile of its null: AI-executed tasks form longer contiguous chains than random reordering would produce, so the chaining pattern is not a compositional artifact.
The result survives pruning.
The observed chain length stays in the upper tail of its null across the inclusive and moderate cuts, at the 87th to 100th percentiles under the Daily$+$ logic and the 58th to 100th under SeveralDaily$+$, and falls back toward the middle of the null only in the Hourly$+$ cuts (the 40th to 68th percentiles), which are the smallest samples.

\begin{figure}[h!]
\caption{Average AI Chain Length Versus Its Task Position-Reshuffle Null Across Frequency Cuts}
\label{fig:chainlength_frequency}
\begin{center}
\includegraphics[width=0.82\textwidth]{plots/fragmentationIndex_weeklyTasks/chainLength_placebo_forest.png}
\end{center}
\footnotesize{
\emph{Notes:} Each row is a frequency cut, with the ``All tasks'' baseline from the main manuscript on top.
The left labels report the number of occupations in each cut.
The dot is the observed average AI chain length for that cut: the mean length of a maximal run of contiguous AI-executed (augmented or automated) steps.
The grey bar spans the 10th to 90th percentiles of the within-occupation task position-reshuffle null ($1{,}000$ draws), and the vertical tick marks the null mean.
The dot is red if it falls outside the null band and blue if inside.
The right margin reports the observed value and its percentile in the null.
}
\end{figure}

\subsection{Prediction \#2: The Empirical Fragmentation Index}
\label{app:frequency_robustness_frag}

Figure~\ref{fig:frag_frequency} summarizes the fragmentation regressions as a heatmap over the logic-by-threshold grid, with the unpruned ``all tasks'' specification from the main analysis in the top row, so that the effect of progressively dropping infrequent tasks can be read down the rows, which impose a more demanding frequency logic, and across the columns, which raise the threshold.
Each cell reports the standardized coefficient on the empirical fragmentation index from re-estimating regression~\eqref{eq:fragmentation_index_regression} for the corresponding cut, with significance stars and the number of occupations printed beneath.
We use our preferred exposure-based measure throughout (Definition~1 of Section~\ref{sec:fragmentation_prediction}, under which both E1- and E2-exposed tasks may form AI chains), and the three heatmaps correspond to the no-fixed-effects, SOC major-group, and SOC minor-group specifications.
A blue cell denotes a negative coefficient, the sign predicted by the model and found in the full-sample estimates, and a red cell a positive one; darker shading indicates a larger magnitude.
The top all-tasks row reproduces the full-sample standardized estimates of $-0.38$, $-0.32$, and $-0.23$, all significant at the 1\% level, here estimated on the $871$ occupations that retain at least five tasks (and matching the headline estimates of Table~\ref{tab:fragmentation_index_regression_exposure}).
In the pruned samples the coefficient remains negative in all but the two smallest cells (Hourly$+$ at the 65\% threshold under SOC fixed effects, $+0.02$ and $+0.01$ with 20 occupations) and stays statistically significant across the inclusive and moderate cuts.
Significance fades only in the sparsest corners, the Hourly$+$ cuts at the 50\% and 65\% thresholds and SeveralDaily$+$ at 65\% (75, 20, and 112 occupations), and more broadly under minor-group fixed effects, where the within-family variation is thinnest.
The fragmentation relationship therefore survives the restriction to frequent tasks.

\begin{figure}[h!]
\caption{Empirical Fragmentation Index and AI Execution Across Frequency Cuts}
\label{fig:frag_frequency}
\begin{center}
\includegraphics[width=\textwidth]{plots/fragmentationIndex_weeklyTasks/frag_logic_threshold_heatmap_def1.png}
\end{center}
\footnotesize{
\emph{Notes:} Each cell reports the standardized coefficient on the empirical fragmentation index from re-estimating regression~\eqref{eq:fragmentation_index_regression} on the corresponding frequency cut.
The index is Definition~1, the preferred exposure-based measure.
The number of occupations appears beneath each estimate, with significance stars $^{*}p{<}.1$, $^{**}p{<}.05$, $^{***}p{<}.01$.
The three heatmaps report the no-fixed-effects, SOC major-group, and SOC minor-group specifications.
Within each heatmap, rows correspond to the three frequency logics and columns to the four thresholds.
The top ``All tasks'' row is the unpruned specification of Section~\ref{sec:fragmentation_prediction}.
Blue denotes a negative coefficient (the predicted sign) and red a positive one.
}
\end{figure}

\subsection{Prediction \#3: The Immediate-Neighbor Effect}
\label{app:frequency_robustness_neighbor}

Figure~\ref{fig:neighbor_frequency_heatmap} repeats the exercise for the neighbor prediction.
Each cell reports the average marginal effect, from re-estimating regression~\eqref{eq:DWA_regression_ai}, of having the immediately preceding ($t{-}1$, Panel~(A)) or following ($t{+}1$, Panel~(B)) task be AI-executed on the focal task's probability of AI execution.
A red cell denotes a positive effect, the sign found in the full sample.
The top row reproduces the full-sample immediate-neighbor effects of $+0.12$ under no fixed effects, which attenuate to between $+0.04$ and $+0.06$ once SOC fixed effects absorb cross-family differences.
The positive immediate-neighbor effect persists across the inclusive and moderate cuts and often grows in magnitude as the workflow is restricted to frequent tasks; under SeveralDaily$+$ $\geq 50\%$, for instance, the previous-task effect rises to $+0.14$ in the no-fixed-effects specification.
The estimates weaken and lose significance only in the sparsest cells, where the number of surviving task observations falls into the low hundreds.
The effects of farther neighbors two positions away remain small and, under fixed effects, are statistically indistinguishable from zero across the grid, in line with the main-text finding that only immediate neighbors matter and with the short average AI chain length of $1.45$.\footnote{ To save space, we do not report the results on steps two position away from the focal step and only focus on the outcomes for immediate neighbors.}

\begin{figure}[p]
\caption{Neighbor AI Execution and Focal-Task AI Execution Across Frequency Cuts}
\label{fig:neighbor_frequency_heatmap}
\begin{center}
\begin{subfigure}[b]{\textwidth}
    \centering
    \captionsetup{labelformat=empty}
    \caption{Panel (A): Previous Task ($t{-}1$)}
    \includegraphics[width=\textwidth]{plots/execTypeVaryingDWA_weeklyTasks/neighbor_logic_threshold_heatmap_prev_bySpec.png}
\end{subfigure}

\vspace{0.8em}

\begin{subfigure}[b]{\textwidth}
    \centering
    \captionsetup{labelformat=empty}
    \caption{Panel (B): Next Task ($t{+}1$)}
    \includegraphics[width=\textwidth]{plots/execTypeVaryingDWA_weeklyTasks/neighbor_logic_threshold_heatmap_next_bySpec.png}
\end{subfigure}
\end{center}
\footnotesize{
\emph{Notes:} Each cell reports the average marginal effect of having the immediately previous task ($t{-}1$, Panel~(A)) or next task ($t{+}1$, Panel~(B)) be AI-executed on the focal task's probability of AI execution, from re-estimating regression~\eqref{eq:DWA_regression_ai} on the corresponding frequency cut.
The number of task observations appears beneath each estimate, with significance stars $^{*}p{<}.1$, $^{**}p{<}.05$, $^{***}p{<}.01$.
Standard errors are clustered at the DWA level.
Within each panel, the three heatmaps report the no-fixed-effects, SOC major-group, and SOC minor-group specifications.
Rows correspond to the three frequency logics and columns to the four thresholds.
The top ``All tasks'' row is the unpruned specification of Section~\ref{sec:DWA_prediction}.
Red denotes a positive effect (the predicted sign) and blue a negative one.
A ``$-$'' marks a cut with too few observations to estimate.
}
\end{figure}

Figure~\ref{fig:neighbor_placebo_forest} benchmarks the previous-task ($t{-}1$, Panel~(A)) and next-task ($t{+}1$, Panel~(B)) effects against their position-reshuffle nulls, one column per specification.
Both effects sit far in the upper tail of their nulls in the full sample, at the 100th percentile in all three specifications, and stay there across the inclusive and moderate cuts.
The previous-task effect lies outside its 10-to-90 null across all four Daily$+$ cuts (the 91st to 99th percentiles) and at SeveralDaily$+$ $\geq 20\%$ and $\geq 50\%$ (the 93rd to 98th), and the next-task effect behaves similarly, escaping its null across the Daily$+$ cuts.
The effects retreat toward the middle of their nulls only at the sparsest cuts, the Hourly$+$ $\geq 35\%$ and $\geq 50\%$ cuts and SeveralDaily$+$ $\geq 35\%$, where the surviving task observations number in the low hundreds and the reshuffle bands are correspondingly wide.
The positive effect of immediate-neighbor AI execution on AI execution likelihood of the focal step persists across the inclusive and moderate cuts and weakens only through lost statistical power at the strictest thresholds, without reversing sign.

In short, all three implications of our model appears to operate among the subset of frequently-performed tasks, and does not appear to be an artifact of the rarely-performed activities catalogued in O*NET.

\begin{figure}[p]
\caption{Immediate-Neighbor Effects: Observed Marginal Effects Versus Their Position-Reshuffle Nulls Across Frequency Cuts}
\label{fig:neighbor_placebo_forest}
\begin{center}
\begin{subfigure}[b]{\textwidth}
    \centering
    \captionsetup{labelformat=empty}
    \caption{Panel (A): Previous Task ($t{-}1$)}
    \includegraphics[width=\textwidth]{plots/placebo_summary/placebo_summary_forest_neighbor_t1_byCut.png}
\end{subfigure}

\vspace{0.8em}

\begin{subfigure}[b]{\textwidth}
    \centering
    \captionsetup{labelformat=empty}
    \caption{Panel (B): Next Task ($t{+}1$)}
    \includegraphics[width=\textwidth]{plots/placebo_summary/placebo_summary_forest_neighbor_t2_byCut.png}
\end{subfigure}
\end{center}
\footnotesize{
\emph{Notes:} Panel~(A) reports the average marginal effect of having the previous task ($t{-}1$) be AI-executed, and Panel~(B) the corresponding effect for the next task ($t{+}1$).
In each panel, rows are frequency cuts, with the ``All tasks'' baseline from the main manuscript on top.
The three columns are the no-fixed-effects, SOC major-group, and SOC minor-group specifications.
The grey bar spans the 10th to 90th percentiles of the within-occupation position-reshuffle null ($1{,}000$ draws), and the vertical tick marks the null mean.
The dot is the observed average marginal effect, red if outside the band and blue if inside.
The left labels report the number of occupations in each cut, and the right margin reports the observed percentile in the null.
A blank row denotes a cut with too few neighbored task observations to estimate.
}
\end{figure}
